\documentclass[12pt]{article}
\usepackage[utf8]{inputenc}
\usepackage[T1]{fontenc}
\usepackage{hyperref}
\title{Taikat Language Resources for K-12 Curriculum Development}
\author{Compiled by Claude Sonnet 5.5 (claude-sonnet-5-5) for Andr\'as Kornai}
\date{}
\begin{document}
\maketitle

\section{Introduction}
Taikat (also Arso; ISO 639-3 aos) is a Border-family Papuan language of Papua Province, Indonesia, spoken in the Arso and Arso Barat districts of Keerom Regency. Wikipedia cites about 500 speakers (2000); Joshua Project classes it as endangered with youth shifting to other languages. Very little is available online.

\section{Primary Educational Resources}
None found.

\section{Linguistic and Reference Materials}
\begin{itemize}
\item Wikipedia, Taikat language (notes six vowels): \url{https://en.wikipedia.org/wiki/Taikat_language}
\item Wikipedia, Border languages (family classification): \url{https://en.wikipedia.org/wiki/Border_languages_(New_Guinea)}
\end{itemize}

\section{Religious and Community Resources}
Joshua Project people-group profile: \url{https://joshuaproject.net/people_groups/15196}

\section{Digital and Audio Resources}
None found.

\section{Curriculum Development Context}
No orthography, primer, dictionary or school programme was found in the search results. Any curriculum would start from fresh documentation with the community.
\end{document}